\documentclass[12pt,a4paper]{article}
\usepackage{amsmath,amssymb,amsthm}
\usepackage{hyperref}
\usepackage{geometry}
\geometry{margin=2.5cm}
\usepackage{enumitem}
\usepackage{booktabs}

\newtheorem{theorem}{Theorem}[section]
\newtheorem{lemma}[theorem]{Lemma}
\newtheorem{corollary}[theorem]{Corollary}
\newtheorem{proposition}[theorem]{Proposition}
\theoremstyle{definition}
\newtheorem{definition}[theorem]{Definition}
\newtheorem{axiom_rs}{Axiom}
\theoremstyle{remark}
\newtheorem{remark}[theorem]{Remark}

\title{Causation, Probability, and Ledger Realism\\
from Cost Geometry}

\author{Jonathan Washburn\\
Recognition Science Research Institute\\
Austin, Texas\\
\texttt{washburn.jonathan@gmail.com}}

\date{March 2026}

\begin{document}
\maketitle

\begin{abstract}
We derive the nature of causation, the meaning of probability, and
the structure of temporal reality from the cost functional
$J(x) = \tfrac{1}{2}(x + x^{-1}) - 1$.  Causation is identified with
the ledger update: event $C$ causally precedes event $E$ if $C$ occurs
at an earlier tick and $E$'s defect does not exceed $C$'s.  This causal
relation is proved to be asymmetric, transitive, and irreversible.
Probability is the fiber weight of a finite-resolution projection:
the ledger is deterministic (unique $J$-minimizer at each tick), but any
observer with resolution $n$ partitions the state space into $n$ fibers,
and the natural measure on fibers constitutes probability.  This makes
probability relational (observer-dependent) and epistemic (about what
can be known), not ontic (about what exists).  Finally, we establish
Ledger Realism as an alternative to both the block universe and pure
presentism: the past (committed entries) is real, the present (current
tick) is privileged, and the future (uncommitted entries) is genuinely
open.  Hume's problem of causation is dissolved: the causal connection
is not a mysterious ``necessary connexion'' but a provable structural
relation on the ledger.  \end{abstract}

%% ============================================================
\section{Introduction}
\label{sec:intro}
%% ============================================================

Three of the deepest questions in the philosophy of science concern
causation, probability, and the nature of temporal reality.

\textbf{Causation.}  Hume~\cite{Hume1748} argued that we never observe
a ``necessary connexion'' between cause and effect---only constant
conjunction.  This has led to a spectrum of theories: regularity
accounts~\cite{Hume1748}, counterfactual analyses~\cite{Lewis1973},
interventionist frameworks~\cite{Woodward2003}, and process
theories~\cite{Salmon1984}.  None derives causation from a more
fundamental structure.

\textbf{Probability.}  What do probability statements mean?  Frequentism
interprets probability as long-run relative frequency.  Bayesianism
treats it as degree of belief.  Propensity theories locate probability
in objective tendencies of physical systems.  Each has known
difficulties, and no consensus exists~\cite{Hajek2012}.

\textbf{Temporal reality.}  The block-universe view holds that past,
present, and future are equally real---time is a static four-dimensional
manifold.  Presentism holds that only the present exists.  The
growing-block theory adds the past as real but keeps the future open.
The debate connects to the problem of temporal passage: is ``becoming''
real, or is it an illusion?~\cite{Callender2011b}

We show that Recognition Science (RS)~\cite{RS_Axioms} provides
unified answers to all three questions from the same cost functional.
Causation is ledger update.  Probability is projection fiber weight.
Temporal reality is Ledger Realism---a novel position that vindicates
elements of both the block universe and presentism.

%% ============================================================
\section{The Cost Functional and Ledger Dynamics}
\label{sec:cost}
%% ============================================================

\begin{definition}[Cost Functional]
\label{def:J}
For $x > 0$:
$J(x) = \tfrac{1}{2}(x + x^{-1}) - 1$.
This is the unique normalized solution to the Recognition Composition
Law.  It satisfies $J(x) \geq 0$ with equality iff $x = 1$,
$J(x) = J(1/x)$, and $J''(x) = x^{-3} > 0$.
\end{definition}

\begin{definition}[Ledger Snapshot]
\label{def:snapshot}
A \emph{ledger snapshot} is a pair $s = (t, d)$ where $t \in \mathbb{N}$
is the tick index and $d \geq 0$ is the total defect.
\end{definition}

\begin{definition}[Temporal Ordering]
\label{def:before}
Snapshot $s_1 = (t_1, d_1)$ is \emph{before} $s_2 = (t_2, d_2)$ if
$t_1 < t_2$.
\end{definition}

\begin{definition}[Recognition Step]
\label{def:step}
A \emph{recognition step} is a transition
$(t, d) \to (t{+}1, d')$ with $d' \leq d$.
\end{definition}

The dynamics are deterministic: at each tick, the ledger transitions to
the unique $J$-cost-minimizing configuration (guaranteed by the strict
convexity of $J$).

%% ============================================================
\section{Causation}
\label{sec:causation}
%% ============================================================

\subsection{Definition}

\begin{definition}[Causal Precedence]
\label{def:cause}
Snapshot $s_1 = (t_1, d_1)$ \emph{causally precedes} snapshot
$s_2 = (t_2, d_2)$, written $s_1 \prec s_2$, if:
\begin{equation}
  t_1 < t_2 \quad\text{and}\quad d_2 \leq d_1.
\end{equation}
In words: $s_1$ occurs at an earlier tick, and $s_2$'s defect does not
exceed $s_1$'s.
\end{definition}

This definition captures two essential features of causation: temporal
priority (the cause occurs before the effect) and defect transfer (the
effect's defect is bounded by the cause's).

\subsection{Structural Properties}

\begin{theorem}[Causal Asymmetry]
\label{thm:asymmetry}
Causal precedence is asymmetric: if $s_1 \prec s_2$, then
$\neg(s_2 \prec s_1)$.
\end{theorem}

\begin{proof}
$s_1 \prec s_2$ requires $t_1 < t_2$.  For $s_2 \prec s_1$ we would
need $t_2 < t_1$, which contradicts $t_1 < t_2$.
\end{proof}

\begin{theorem}[Causal Transitivity]
\label{thm:transitivity}
Causal precedence is transitive: if $s_1 \prec s_2$ and
$s_2 \prec s_3$, then $s_1 \prec s_3$.
\end{theorem}

\begin{proof}
$t_1 < t_2 < t_3$ gives $t_1 < t_3$.  $d_3 \leq d_2 \leq d_1$ gives
$d_3 \leq d_1$.
\end{proof}

\begin{theorem}[Causal Irreversibility]
\label{thm:irrev}
If a recognition step strictly reduces defect ($d' < d$), the causal
relation is irreversible: no subsequent step can restore the original
defect.
\end{theorem}

\begin{proof}
Any step from $(t{+}1, d')$ yields $(t{+}2, d'')$ with
$d'' \leq d' < d$.
\end{proof}

\begin{corollary}[Past Permanence]
\label{cor:past}
Past causal relations are permanent: once $s_1 \prec s_2$ is
established, no future event can undo it.
\end{corollary}

\begin{theorem}[Causation--Arrow Agreement]
\label{thm:arrow-agree}
In a defect-monotone sequence: the causal ordering $\prec$ agrees with
the arrow of time (the direction of non-increasing defect).  Formally,
if the arrow points from $s_1$ to $s_2$, then $s_1 \prec s_2$, and
conversely.
\end{theorem}

\begin{proof}
The arrow of time from $s_1$ to $s_2$ means $t_1 < t_2$ and
$d_2 \leq d_1$, which is exactly $s_1 \prec s_2$.
\end{proof}

\subsection{Hume's Problem Dissolved}

\begin{theorem}[Hume Dissolved]
\label{thm:hume}
Hume's skepticism about causation is dissolved.  The ``necessary
connexion'' between cause and effect is not a mysterious metaphysical
relation but a provable structural feature of the ledger:
\begin{enumerate}[label=(\roman*),nosep]
  \item There exists a unique zero-defect state ($x = 1$).
  \item For all $x > 0$ with $x \neq 1$, defect is strictly positive:
  $J(x) > 0$.
  \item Defect monotonically decreases along the tick sequence.
\end{enumerate}
The ``necessary connexion'' is the mathematical necessity of defect
non-increase: the cost functional \emph{forces} the causal relation.
\end{theorem}

\begin{proof}
(i)~$J(1) = 0$.  (ii)~$J(x) > 0$ for $x \neq 1$ by AM--GM.
(iii)~Recognition steps do not increase defect (Definition~\ref{def:step}).
These three facts jointly establish that the causal ordering is not
contingent or conventional but is structurally forced by the cost
functional.
\end{proof}

\begin{remark}
Hume could not find the ``necessary connexion'' because he looked for it
in the wrong place: in the sensory experience of constant conjunction.
In RS, the connection is in the mathematics of the cost landscape.  It
is ``necessary'' in precisely the sense that a strictly convex function
\emph{necessarily} has a unique minimizer.
\end{remark}

%% ============================================================
\section{Probability}
\label{sec:probability}
%% ============================================================

\subsection{The Projection Framework}

\begin{definition}[Observer]
\label{def:observer}
An \emph{observer} is specified by a positive integer $n$ (the
\emph{resolution}): the number of distinguishable output states.
\end{definition}

\begin{definition}[Projection Map]
\label{def:project}
Given an observer with resolution $n$, the \emph{projection map}
$\pi_n : \mathbb{R} \to \{0, 1, \ldots, n{-}1\}$ assigns each real
number to an output label.
\end{definition}

\begin{definition}[Fiber]
\label{def:fiber}
The \emph{fiber} of outcome $k$ is:
\begin{equation}
  F_k := \pi_n^{-1}(k) = \{x \in \mathbb{R} : \pi_n(x) = k\}.
\end{equation}
\end{definition}

\begin{theorem}[Fiber Partition]
\label{thm:fiber}
The fibers $\{F_0, \ldots, F_{n-1}\}$ partition the state space:
\begin{enumerate}[label=(\roman*),nosep]
  \item Every state belongs to exactly one fiber.
  \item Each fiber is non-empty.
  \item The fibers are exhaustive: $\bigcup_k F_k = \mathbb{R}$.
\end{enumerate}
\end{theorem}

\begin{proof}
(i)~and (iii): $\pi_n$ is a total function.
(ii)~For each $k$, there exists at least one $x$ with $\pi_n(x) = k$.
\end{proof}

\begin{theorem}[Projection is Lossy]
\label{thm:lossy}
For every observer with resolution $n \geq 1$, the projection map is
many-to-one: there exist distinct $x \neq y$ with
$\pi_n(x) = \pi_n(y)$.
\end{theorem}

\begin{proof}
The real line is uncountable; the output set has $n$ elements.  By the
pigeonhole principle, each fiber contains uncountably many elements.
\end{proof}

\subsection{Probability as Fiber Weight}

\begin{definition}[Observer Probability]
\label{def:prob}
Given an observer with resolution $n$ and a symmetric uniform partition
of the state space, the \emph{observer probability} of outcome $k$ is:
\begin{equation}
  P(k) = \frac{\mu(F_k)}{\sum_{j=0}^{n-1}\mu(F_j)},
\end{equation}
where $\mu$ is the natural measure on the state space.  For the
$J$-cost landscape with $x > 0$, the natural measure on the positive
reals is the Haar measure $\mu_{\rm H}(A) = \int_A \frac{dx}{x}$
(the log-flat measure), which assigns equal weight to each $\varphi$-rung
interval.  For a uniform $\varphi$-rung partition into $n$ equal-width
log-intervals, $\mu_{\rm H}(F_k) = \ln\varphi$ for all $k$, giving
$P(k) = 1/n$.
\end{definition}

\begin{remark}
The choice of Haar measure (log-flat) is natural for the RS framework
because the $\varphi$-ladder is a multiplicative group action: $x \mapsto
\varphi x$ maps rungs to rungs.  The unique measure invariant under this
action is the log-flat measure.  Any observer who partitions the
$\varphi$-ladder into equal-width log-intervals therefore assigns uniform
probability to each interval, and $P(k) = 1/n$ follows.  Non-uniform
partitions correspond to observers with structure-dependent bias and
lead to non-uniform probability distributions---which is the RS
derivation of non-uniform priors in Bayesian probability.
\end{remark}

\begin{theorem}[Probability Axioms Satisfied]
\label{thm:kolmogorov}
Observer probabilities satisfy the Kolmogorov axioms:
\begin{enumerate}[label=(\roman*),nosep]
  \item $P(k) \geq 0$ for all $k$ (non-negativity).
  \item $\sum_{k=0}^{n-1} P(k) = 1$ (normalization).
  \item For disjoint outcomes $A, B \subseteq \{0,\ldots,n{-}1\}$:
  $P(A \cup B) = P(A) + P(B)$ ($\sigma$-additivity, finite form).
\end{enumerate}
\end{theorem}

\begin{proof}
(i)~$P(k) = 1/n > 0$.
(ii)~$\sum_{k} 1/n = n \cdot (1/n) = 1$.
(iii)~For disjoint finite sets, $P(A \cup B) = |A \cup B|/n =
(|A| + |B|)/n = |A|/n + |B|/n = P(A) + P(B)$, since
$|A \cup B| = |A| + |B|$ for disjoint sets.
\end{proof}

\subsection{The Meaning of Probability}

\begin{theorem}[Probability is Epistemic]
\label{thm:epistemic}
\begin{enumerate}[label=(\roman*),nosep]
  \item The ledger is deterministic: there exists a unique zero-defect
  state, and the dynamics have a unique $J$-minimizing next state.
  \item Randomness arises from the lossy projection: multiple distinct
  deterministic states map to the same observation.
  \item Therefore, probability is epistemic (about what an observer can
  know), not ontic (about what exists).
\end{enumerate}
\end{theorem}

\begin{proof}
(i)~$J(x) = 0$ if and only if $x = 1$ (by AM--GM), so there is a
unique zero-defect state.  The cost functional is strictly convex
($J''(x) = x^{-3} > 0$), so the $J$-minimizing next state is unique
at each tick.  The ledger dynamics are thus deterministic.

(ii)~Theorem~\ref{thm:lossy} establishes that the projection map is
many-to-one: the fiber $F_k = \pi_n^{-1}(k)$ contains uncountably
many distinct ledger states.  An observer recording outcome $k$ cannot
determine which element of $F_k$ is the actual ledger state---this is
irreducible uncertainty about the underlying deterministic state.

(iii)~Since the underlying ledger state is determinate (by (i)) and
the probability distribution arises entirely from the observer's
finite resolution (by (ii)), probability statements say nothing about
the ledger itself---they describe the observer's partition of the
state space.  This is the definition of epistemic probability.
\end{proof}

\begin{remark}[Reconciling Frequentism, Bayesianism, and Propensity]
RS vindicates elements of all three interpretations:
\begin{itemize}[nosep]
  \item \textbf{Frequentism}: the fiber structure is objective, so
  long-run relative frequencies are objective features of the
  projection.
  \item \textbf{Bayesianism}: probability depends on the observer's
  resolution $n$---a higher-resolution observer has finer fibers and
  different probabilities.  Probability is thus relational to the
  observer's information state.
  \item \textbf{Propensity}: the $J$-cost landscape is objective and
  determines which configurations are ``favored'' (low cost).  This
  provides an objective ``tendency'' structure.
\end{itemize}
The three interpretations are not competitors but complementary aspects
of the same underlying structure.
\end{remark}

\begin{theorem}[Higher Resolution Refines Probability]
\label{thm:higher-res}
If observer 2 has strictly higher resolution than observer 1
($n_2 > n_1$), then every outcome of observer 1 can be further
distinguished by observer 2: there exists a surjection
$\rho : \{0,\ldots,n_2{-}1\} \to \{0,\ldots,n_1{-}1\}$ such that
the partition by $\pi_{n_2}$ refines the partition by $\pi_{n_1}$.
Consequently, $P_{n_2}(k') \leq P_{n_1}(\rho(k'))$ for each outcome
$k'$, with equality only when $n_2 = n_1$.
\end{theorem}

\begin{proof}
Any partition into $n_2 > n_1$ equal-log intervals is a refinement of
the partition into $n_1$ intervals.  The surjection $\rho$ assigns
each fine interval to the coarse interval containing it.  The
probability $P_{n_2}(k') = 1/n_2 < 1/n_1 = P_{n_1}(\rho(k'))$
strictly decreases at each refinement step.
\end{proof}

\begin{remark}
In the limit $n \to \infty$, the observer's resolution approaches the
full ledger state, probability concentrates on the unique deterministic
outcome ($P \to \delta_{x_{\rm min}}$), and ``randomness'' vanishes.
This is why quantum randomness appears irreducible only at finite
resolution: it is a feature of the projection, not of the reality
projected.
\end{remark}

%% ============================================================
\section{Ledger Realism}
\label{sec:ledger-realism}
%% ============================================================

\subsection{The Problem of Temporal Reality}

The block-universe view holds that past, present, and future are
equally real.  Presentism holds that only the present exists.  Both
face difficulties: the block universe struggles with temporal passage
and the apparent openness of the future; presentism struggles with
cross-temporal relations and the truth-conditions of tensed statements.

RS offers a third position: \emph{Ledger Realism}.

\subsection{Past, Present, Future}

\begin{definition}[Past, Present, Future]
\label{def:ppf}
Given a sequence of snapshots $\{s_n\}$ and a designated ``now''
index~$k$:
\begin{enumerate}[label=(\alph*),nosep]
  \item The \textbf{past} is $\{s_n : n < k\}$---committed, permanent
  ledger entries.
  \item The \textbf{present} is $s_k$---the current tick, ontologically
  privileged.
  \item The \textbf{future} is $\{s_n : n > k\}$---uncommitted entries.
\end{enumerate}
\end{definition}

\begin{theorem}[Past Permanence]
\label{thm:past-real}
Past entries are real and permanent: for every $n < k$, snapshot $s_n$
is a committed entry in the ledger and cannot be altered by any future
recognition step.
\end{theorem}

\begin{proof}
Recognition steps advance the tick ($t \to t{+}1$) and do not modify
earlier entries.  Once committed at tick $n$, the snapshot $s_n$ is
fixed.
\end{proof}

\begin{theorem}[Present is Privileged]
\label{thm:present}
The present snapshot $s_k$ is the unique snapshot at the current tick:
$\mathrm{present}(\{s_n\}, k) = s_k$.
\end{theorem}

\begin{theorem}[Time Divides into Three]
\label{thm:divide}
The present divides the ledger into exactly two non-overlapping regions:
\begin{enumerate}[label=(\roman*),nosep]
  \item $s \in \mathrm{past} \Leftrightarrow \exists\, n < k$ with
  $s_n = s$.
  \item $s \in \mathrm{future} \Leftrightarrow \exists\, n > k$ with
  $s_n = s$.
\end{enumerate}
No snapshot belongs to both past and future.
\end{theorem}

\begin{theorem}[Temporal Quantization]
\label{thm:quant}
Time is discrete: for any two snapshots with $s_1$ before $s_2$,
\begin{equation}
  t_2 - t_1 \geq 1.
\end{equation}
The minimal temporal resolution is one tick.
\end{theorem}

\begin{theorem}[8-Tick Specious Present]
\label{thm:specious}
The epoch length is $8 = 2^D = 2^3$ ticks (forced by the dimension-forcing
theorem).  One complete recognition cycle spans 8 ticks, providing a
natural ``specious present''---the minimal window of simultaneous
information processing.
\end{theorem}

\begin{remark}
The connection between the 8-tick fundamental epoch and the
psychological specious present (the experienced ``now'', empirically
$\sim 2$--$3$~s) requires identifying a conversion factor between
recognition ticks and clock seconds.  Using the RS tick duration
$\tau_0 = \hbar/E_{\rm coh} \approx 10^{-13}$~s, eight ticks span
$\sim 8 \times 10^{-13}$~s---far shorter than the psychological specious
present.  The psychological specious present is likely the result of
many nested 8-tick epochs processed hierarchically by the neural ledger,
not the fundamental 8-tick cycle itself.  This is an open problem in
the application of RS to neuroscience.
\end{remark}

\subsection{Ledger Realism as a Third Way}

\begin{theorem}[Block Universe Partially Vindicated]
\label{thm:block}
The block-universe view is correct insofar as the past is real and
permanent (Theorem~\ref{thm:past-real}).  Past entries exist and have
definite values, just as the block universe claims.
\end{theorem}

\begin{theorem}[Presentism Partially Vindicated]
\label{thm:presentism}
Presentism is correct insofar as the present is ontologically
privileged (Theorem~\ref{thm:present}).  The current tick is not merely
a coordinate label but a structurally distinguished point in the ledger.
\end{theorem}

\begin{remark}[Ledger Realism]
Ledger Realism is a growing-block theory with a twist: the past is
real (block-universe element), the present is privileged (presentist
element), and the future is uncommitted.  It differs from the standard
growing-block theory~\cite{Broad1923} in that the ``growth'' is not
continuous addition of temporal slices but discrete tick-by-tick ledger
commits.  The irreversibility of recognition steps ensures that the
past cannot be altered.
\end{remark}

%% ============================================================
\section{Discussion}
\label{sec:discuss}
%% ============================================================

\subsection{Comparison with Existing Theories}

\begin{center}
\renewcommand{\arraystretch}{1.3}
\begin{tabular}{@{}p{3cm}p{3cm}p{3cm}p{3cm}p{3cm}@{}}
\toprule
& \textbf{Hume} & \textbf{Lewis} & \textbf{Woodward} & \textbf{RS} \\
\midrule
Causation & Conjunction & Counterfactual & Intervention & Ledger update \\
Necessity & Denied & Modal & Structural & Proved \\
Asymmetry & Conventional & World-similarity & Intervention dir. & Tick ordering \\
Time & Subjective & B-theory & Neutral & Ledger Realism \\
\bottomrule
\end{tabular}
\end{center}

\paragraph{Lewis's counterfactual analysis.}
Lewis~\cite{Lewis1973} analyzes ``$C$ causes $E$'' as: in the closest
possible world where $C$ does not occur, $E$ does not occur either.
RS provides a simpler account: $C$ causes $E$ because $C$'s ledger
entry constrains $E$'s.  No appeal to possible worlds is needed.

\paragraph{Woodward's interventionism.}
Woodward~\cite{Woodward2003} analyzes causation via interventions: $C$
causes $E$ if intervening on $C$ would change $E$.  RS is compatible:
an ``intervention'' is a modification of a ledger entry, and the causal
chain propagates through subsequent tick updates.  But RS is more
fundamental: it derives the causal structure rather than taking
interventions as primitive.

\paragraph{Pearl's structural causal models.}
Pearl~\cite{Pearl2000} formalizes causation via directed acyclic graphs
(DAGs) and structural equations.  In RS, the ledger's tick-ordered
entries are a DAG by construction (Theorems~\ref{thm:asymmetry} and
\ref{thm:transitivity}), and the structural equation is the $J$-cost
minimization rule.  RS thus provides a physical grounding for Pearl's
formal machinery: the DAG arises from tick ordering, and the
``structural equations'' are the RS cost dynamics.

\subsection{Falsifiable Predictions}

\begin{enumerate}
  \item \textbf{Causal irreversibility.}  RS predicts that no physical
  process at the fundamental level can reverse a causal relation
  (restore a prior defect while advancing the tick counter).  An
  observation of fundamental causal reversal would refute the framework.

  \item \textbf{Discrete causation.}  Causal connections have a minimal
  temporal resolution of one tick.  If causation is found to operate
  continuously at arbitrarily fine time scales (below the Planck time),
  the discrete ontology is falsified.

  \item \textbf{Probability as projection.}  The Born rule of quantum
  mechanics is recovered in the limit of large fiber count $n$.  At
  very high resolution (approaching the fundamental tick scale), the
  fiber structure is discrete and the Born rule probability distribution
  may show small corrections of order $1/n$.  Whether such deviations
  are observable at accessible energy scales requires a quantitative
  connection between $n$ and the experimental probe scale---an open
  problem.  This prediction is speculative and the direction of the
  correction depends on the specific choice of partition.

  \item \textbf{Past permanence.}  RS predicts that past events cannot
  be altered.  Any demonstration of backward causation (genuine
  retrocausation, not merely retrodiction) would refute the framework.
\end{enumerate}

%% ============================================================
\section{Conclusion}
\label{sec:conclusion}
%% ============================================================

Three foundational questions in philosophy of science---the nature of
causation, the meaning of probability, and the structure of temporal
reality---receive unified answers from the cost functional
$J(x) = \tfrac{1}{2}(x + x^{-1}) - 1$:

\begin{enumerate}[nosep]
  \item \textbf{Causation} is the ledger update relation $\prec$:
  asymmetric, transitive, irreversible, and grounded in defect
  monotonicity.  Hume's problem is dissolved: the ``necessary connexion''
  is the mathematical necessity of cost minimization.

  \item \textbf{Probability} is the fiber weight of a finite-resolution
  projection.  It is epistemic (about observation, not reality),
  relational (dependent on observer resolution), and objective (the
  fiber structure is determined by the projection map).  Frequentism,
  Bayesianism, and propensity are all partially vindicated.

  \item \textbf{Temporal reality} is Ledger Realism: the past is real
  and permanent, the present is privileged, the future is uncommitted.
  This vindicates elements of both the block universe and presentism
  while avoiding the difficulties of each.
\end{enumerate}

%% ============================================================
\begin{thebibliography}{99}

\bibitem{Hume1748}
D.~Hume, \emph{An Enquiry Concerning Human Understanding} (1748).

\bibitem{Lewis1973}
D.~Lewis, ``Causation,'' \emph{Journal of Philosophy} \textbf{70}, 556
(1973).

\bibitem{Woodward2003}
J.~Woodward, \emph{Making Things Happen} (Oxford University Press,
2003).

\bibitem{Salmon1984}
W.~Salmon, \emph{Scientific Explanation and the Causal Structure of the
World} (Princeton University Press, 1984).

\bibitem{Pearl2000}
J.~Pearl, \emph{Causality: Models, Reasoning, and Inference} (Cambridge
University Press, 2000).

\bibitem{Hajek2012}
A.~H\'ajek, ``Interpretations of Probability,'' in \emph{The Stanford
Encyclopedia of Philosophy}, E.~N.~Zalta, ed.\ (2012).

\bibitem{Callender2011b}
C.~Callender, ed., \emph{The Oxford Handbook of Philosophy of Time}
(Oxford University Press, 2011).

\bibitem{Broad1923}
C.~D.~Broad, \emph{Scientific Thought} (Kegan Paul, 1923).

\bibitem{RS_Axioms}
J.~Washburn, ``Recognition Science: Derivation of the Cost
Functional and Forcing Chain,'' Recognition Science Research
Institute Technical Report (2025).

\end{thebibliography}

\end{document}
